\documentclass[10pt,a4paper]{amsart}
\usepackage{amsmath,amssymb,amsthm,graphicx}
\usepackage{verbatim}
\newenvironment{reference}{\comment}{\endcomment}
%\newenvironment{reference}{}{}
\newenvironment{slogan}{\comment}{\endcomment}
\newenvironment{history}{\comment}{\endcomment}

% For commutative diagrams we use Xy-pic
\usepackage[all]{xy}

% We use 2cell for 2-commutative diagrams.
\xyoption{2cell}
% Theorem environments.
%
\theoremstyle{plain}
\newtheorem{theorem}[subsection]{Theorem}
\newtheorem{proposition}[subsection]{Proposition}
\newtheorem{lemma}[subsection]{Lemma}

\theoremstyle{definition}
\newtheorem{definition}[subsection]{Definition}
\newtheorem{example}[subsection]{Example}
\newtheorem{exercise}[subsection]{Exercise}
\newtheorem{situation}[subsection]{Situation}

\theoremstyle{remark}
\newtheorem{remark}[subsection]{Remark}
\newtheorem{remarks}[subsection]{Remarks}

\numberwithin{equation}{subsection}

% Macros
%
\def\lim{\mathop{\mathrm{lim}}\nolimits}
\def\colim{\mathop{\mathrm{colim}}\nolimits}
\def\Spec{\mathop{\mathrm{Spec}}}
\def\Hom{\mathop{\mathrm{Hom}}\nolimits}
\def\Ext{\mathop{\mathrm{Ext}}\nolimits}
\def\SheafHom{\mathop{\mathcal{H}\!\mathit{om}}\nolimits}
\def\SheafExt{\mathop{\mathcal{E}\!\mathit{xt}}\nolimits}
\def\Sch{\mathit{Sch}}
\def\Mor{\mathop{\mathrm{Mor}}\nolimits}
\def\Ob{\mathop{\mathrm{Ob}}\nolimits}
\def\Sh{\mathop{\mathit{Sh}}\nolimits}
\def\NL{\mathop{N\!L}\nolimits}
\def\CH{\mathop{\mathrm{CH}}\nolimits}
\def\proetale{{pro\text{-}\acute{e}tale}}
\def\etale{{\acute{e}tale}}
\def\QCoh{\mathit{QCoh}}
\def\Ker{\mathop{\mathrm{Ker}}}
\def\Im{\mathop{\mathrm{Im}}}
\def\Coker{\mathop{\mathrm{Coker}}}
\def\Coim{\mathop{\mathrm{Coim}}}

% Boxtimes
%
\DeclareMathSymbol{\boxtimes}{\mathbin}{AMSa}{"02}

%
% Macros for moduli stacks/spaces
%
\def\QCohstack{\mathcal{QC}\!\mathit{oh}}
\def\Cohstack{\mathcal{C}\!\mathit{oh}}
\def\Spacesstack{\mathcal{S}\!\mathit{paces}}
\def\Quotfunctor{\mathrm{Quot}}
\def\Hilbfunctor{\mathrm{Hilb}}
\def\Curvesstack{\mathcal{C}\!\mathit{urves}}
\def\Polarizedstack{\mathcal{P}\!\mathit{olarized}}
\def\Complexesstack{\mathcal{C}\!\mathit{omplexes}}
% \Pic is the operator that assigns to X its picard group, usage \Pic(X)
% \Picardstack_{X/B} denotes the Picard stack of X over B
% \Picardfunctor_{X/B} denotes the Picard functor of X over B
\def\Pic{\mathop{\mathrm{Pic}}\nolimits}
\def\Picardstack{\mathcal{P}\!\mathit{ic}}
\def\Picardfunctor{\mathrm{Pic}}
\def\Deformationcategory{\mathcal{D}\!\mathit{ef}}
\usepackage{cleveref}
\crefname{originallemma}{original lemma}{original lemmas}
\Crefname{originallemma}{Original lemma}{Original lemmas}
\makeatletter
\newlabel{lemma-ass-support}{{0586}{source}{}{sourcepointer-1}{}}
\newlabel{lemma-ass-support@cref}{{[originallemma][0][]0586}{[1][0][]source}{}{}{}}
\makeatother
\crefname{originallemma}{original lemma}{original lemmas}
\Crefname{originallemma}{Original lemma}{Original lemmas}
\makeatletter
\newlabel{lemma-filter-minimal-primes-in-support}{{00L7}{source}{}{sourcepointer-2}{}}
\newlabel{lemma-filter-minimal-primes-in-support@cref}{{[originallemma][0][]00L7}{[1][0][]source}{}{}{}}
\makeatother
\title{Minimal primes of the support, of the associated primes and of a filtration all agree}
\author{The Stacks Project Authors}
\date{stacks-project:algebra.tex:local-snapshot-20261012; GNU FDL 1.2 or later}
\begin{document}
\maketitle
\noindent\textit{Editorial note.} All delivered mathematical text is the author's own, copied from the pinned Stacks Project chapter file algebra.tex; no statement or proof step was rewritten, reordered or removed.\par\medskip
\noindent\textit{Editorial note.} Preamble provenance, as required for this source: the frozen snapshot stages the chapter's own preamble file beside it (preamble.tex, 217 lines, 8116 bytes, the same shared Stacks-wide preamble the published chow and etale-cohomology cards use); the first line of the pinned chapter file is \textbackslash{}input\{preamble\}, so no preamble text lives inside the chapter itself. The card takes that file's lines 8-18 (the verbatim package and the three \textbackslash{}comment-based environments reference/slogan/history at lines 9-12, section 13.2, together with the Xy-pic setup) and its lines 154 to the runtime end of the file (the native theorem declarations and mathematical macros). Not one character was changed, no package was added, no macro was defined, and the excerpt serves only as preamble, never as a span coordinate.\par\medskip
\noindent\textit{Editorial note.} The same-chapter pointers are located at their \textbackslash{}label lines inside the pinned chapter file; the cross-chapter pointers are located in the project's own tag list staged in the snapshot (sources/stacks-chow/tags.txt), at the verbatim lines carrying their canonical tags.\par\medskip
\noindent\textit{Editorial note.} The counted proof refers to same-chapter statements by label; they are recorded as navigation stubs pointed at their \textbackslash{}label lines, and the arguments stay inside the counted span.\par\medskip
\noindent\textit{Editorial note.} The counted spans contain no \textbackslash{}cite, so the card carries no bibliography.\par\medskip
\noindent\textit{Allowed standard prerequisites (editorial).}\begin{itemize}
\item Associated primes of a filtration with cyclic prime quotients. Invoked in the counted proof at the start (lines 15419-15421). Identifies the minimal primes of the support with those of the primes of a filtration.
\item Minimality of a prime in the support versus in the associated primes. Used in the second half of the counted proof (lines 15423-15447). Transfers minimality between the two sets of primes through annihilators.
\end{itemize}\medskip
\ifdefined\hypertarget\hypertarget{sourcepointer-1}{}\fi
\noindent\textit{Source reference.} Original lemma 0586: sources/stacks-algebra/pinned-original-source.tex, source lines 15335--15335. Navigation pointer for \textbackslash{}ref\{lemma-ass-support\} used inside the counted span (line 15436); the argument that consumes it is retained in the span.\par\medskip
\ifdefined\hypertarget\hypertarget{sourcepointer-2}{}\fi
\noindent\textit{Source reference.} Original lemma 00L7: sources/stacks-algebra/pinned-original-source.tex, source lines 15234--15234. Navigation pointer for \textbackslash{}ref\{lemma-filter-minimal-primes-in-support\} used inside the counted span (line 15420); the argument that consumes it is retained in the span.\par\medskip
\begin{proposition}
\label{proposition-minimal-primes-associated-primes}
Let $R$ be a Noetherian ring.
Let $M$ be a finite $R$-module.
The following sets of primes are the same:
\begin{enumerate}
\item The minimal primes in the support of $M$.
\item The minimal primes in $\text{Ass}(M)$.
\item For any filtration $0 = M_0 \subset M_1 \subset \ldots
\subset M_{n-1} \subset M_n = M$ with $M_i/M_{i-1} \cong R/\mathfrak p_i$
the minimal primes of the set $\{\mathfrak p_i\}$.
\end{enumerate}
\end{proposition}

\begin{proof}
Choose a filtration as in (3).
In Lemma \ref{lemma-filter-minimal-primes-in-support}
we have seen that the sets in (1) and (3) are equal.

\medskip\noindent
Let $\mathfrak p$ be a minimal element of the set $\{\mathfrak p_i\}$.
Let $i$ be minimal such that $\mathfrak p = \mathfrak p_i$.
Pick $m \in M_i$, $m \not \in M_{i-1}$. The annihilator of $m$
is contained in $\mathfrak p_i = \mathfrak p$ and contains
$\mathfrak p_1 \mathfrak p_2 \ldots \mathfrak p_i$. By our choice of
$i$ and $\mathfrak p$ we have $\mathfrak p_j \not \subset \mathfrak p$
for $j < i$ and hence we have
$\mathfrak p_1 \mathfrak p_2 \ldots \mathfrak p_{i - 1}
\not \subset \mathfrak p_i$. Pick
$f \in \mathfrak p_1 \mathfrak p_2 \ldots \mathfrak p_{i - 1}$,
$f \not \in \mathfrak p$. Then $fm$ has annihilator $\mathfrak p$.
In this way we see that $\mathfrak p$ is an associated prime of $M$.
By Lemma \ref{lemma-ass-support} we have $\text{Ass}(M) \subset \text{Supp}(M)$
and hence $\mathfrak p$ is minimal in $\text{Ass}(M)$.
Thus the set of primes in (1) is contained in the set of primes of (2).

\medskip\noindent
Let $\mathfrak p$ be a minimal element of $\text{Ass}(M)$.
Since $\text{Ass}(M) \subset \text{Supp}(M)$ there is a minimal
element $\mathfrak q$ of $\text{Supp}(M)$ with
$\mathfrak q \subset \mathfrak p$. We have just shown that
$\mathfrak q \in \text{Ass}(M)$. Hence $\mathfrak q = \mathfrak p$
by minimality of $\mathfrak p$. Thus the set of primes in (2) is
contained in the set of primes of (1).
\end{proof}

\end{document}
